\documentclass[10pt, letterpaper]{article}

\usepackage[T1]{fontenc}
\usepackage{lmodern}
\usepackage[margin=0.7in, top=0.65in, bottom=0.65in]{geometry}
\usepackage{amsmath, amssymb}
\usepackage{booktabs}
\usepackage{hyperref}
\usepackage{multicol}
\usepackage{fancyvrb}
\usepackage{xcolor}

\setlength{\parskip}{3pt plus 1pt minus 1pt}
\setlength{\parindent}{0pt}
\setlength{\textfloatsep}{4pt}
\setlength{\floatsep}{4pt}
\setlength{\intextsep}{4pt}
\setlength{\abovecaptionskip}{3pt}
\setlength{\belowcaptionskip}{1pt}

% Code block style
\definecolor{codebg}{gray}{0.95}
\DefineVerbatimEnvironment{code}{Verbatim}{fontsize=\fontsize{7}{8.5}\selectfont,
  frame=single, framerule=0.3pt, rulecolor=\color{gray!50},
  fillcolor=\color{codebg}, formatcom=\color{black},
  xleftmargin=0.5em, xrightmargin=0.5em, numbers=left, numbersep=4pt}

\title{\vspace{-2em}\textbf{PML Assignment 3: Logistic Regression vs.\ SVM on Breast Cancer Wisconsin Data}\vspace{-0.5em}}
\author{Brian Butler}
\date{March 2026}

\begin{document}
\maketitle
\vspace{-1em}

%----------------------------------------------------------------------
\section{Overview}
%----------------------------------------------------------------------

We compare \textbf{Logistic Regression}, \textbf{Linear SVM}, and \textbf{RBF-kernel SVM} on the \textbf{Breast Cancer Wisconsin (Diagnostic)} dataset ($n{=}569$ samples, $d{=}30$ features, binary: malignant vs.\ benign).
All models are wrapped in \texttt{Pipeline(StandardScaler, model)} so that feature standardization is performed \emph{inside} cross-validation, preventing data leakage.
Hyperparameters are tuned via \texttt{GridSearchCV} with 5-fold \texttt{StratifiedKFold}, and a single held-out test set (20\%) is evaluated once at the end.

%----------------------------------------------------------------------
\section{Experimental Setup}
%----------------------------------------------------------------------

\subsection{Data Split}
An 80/20 stratified split (\texttt{random\_state=42}) produces 455 training and 114 test samples.
The test set is never used during tuning.

\subsection{Pipeline \& CV Configuration}
\begin{code}
# Pipeline prevents data leakage: scaling fit on each CV fold
pipe = Pipeline([
    ("scaler", StandardScaler()),
    ("clf",    model)   # LR / LinearSVC / SVC
])
# Stratified 5-fold CV ensures class balance in every fold
cv = StratifiedKFold(n_splits=5, shuffle=True, random_state=42)
gs = GridSearchCV(pipe, param_grid, cv=cv, scoring="accuracy")
\end{code}

\subsection{Hyperparameter Grids}

\begin{center}
\small
\textbf{Table 1:} Search grids for each model\\[3pt]
\begin{tabular}{lll}
\toprule
Model & Parameter(s) & Grid Values \\
\midrule
Logistic Regression & $C$ & $\{0.001, 0.01, 0.1, 1, 10, 100\}$ \\
Linear SVC          & $C$ & $\{0.001, 0.01, 0.1, 1, 10, 100\}$ \\
RBF SVC             & $C$ & $\{0.01, 0.1, 1, 10, 100\}$ \\
                    & $\gamma$ & $\{0.001, 0.01, 0.1, 1\}$ \\
\bottomrule
\end{tabular}
\end{center}

\textbf{Note:} \texttt{LinearSVC} lacks \texttt{predict\_proba}, so it is wrapped in \texttt{CalibratedClassifierCV} (3-fold internal CV) to enable ROC--AUC computation.

%----------------------------------------------------------------------
\section{Results}
%----------------------------------------------------------------------

\subsection{Best CV Hyperparameters}

\begin{center}
\small
\textbf{Table 2:} Best hyperparameters selected by 5-fold GridSearchCV\\[3pt]
\begin{tabular}{lcc}
\toprule
Model & Best Params & Best CV Accuracy \\
\midrule
Logistic Regression & $C = 0.1$            & 0.9824 \\
Linear SVC          & $C = 0.01$           & 0.9780 \\
RBF SVC             & $C = 10,\;\gamma = 0.01$ & 0.9758 \\
\bottomrule
\end{tabular}
\end{center}

All three models achieve CV accuracies above 97.5\%, reflecting that the breast cancer dataset is well-separated after standardization.

\subsection{Held-Out Test-Set Performance}

\begin{center}
\small
\textbf{Table 3:} Test-set metrics (114 samples, evaluated once)\\[3pt]
\begin{tabular}{lccc}
\toprule
Model & Test Accuracy & ROC--AUC & Confusion Matrix \\
\midrule
Logistic Regression & 0.9737 & 0.9957 & $\begin{pmatrix}40 & 2\\1 & 71\end{pmatrix}$ \\[6pt]
Linear SVC          & 0.9825 & 0.9960 & $\begin{pmatrix}41 & 1\\1 & 71\end{pmatrix}$ \\[6pt]
RBF SVC             & 0.9825 & 0.9977 & $\begin{pmatrix}41 & 1\\1 & 71\end{pmatrix}$ \\
\bottomrule
\end{tabular}
\end{center}

Confusion matrix format: rows = actual (malignant, benign), columns = predicted.
TN = true negatives (malignant correctly identified); TP = true positives (benign correctly identified).

%----------------------------------------------------------------------
\section{Discussion}
%----------------------------------------------------------------------

\subsection{Model Comparison}

All three classifiers achieve excellent performance ($>97\%$ accuracy, $>99.5\%$ AUC), consistent with existing literature on this dataset.

\textbf{Logistic Regression} achieves the highest CV accuracy (0.9824) but slightly lower test accuracy (0.9737) with 2 false positives versus 1 for the SVMs.
The selected $C{=}0.1$ indicates stronger regularization is preferred, suggesting the decision boundary benefits from shrinkage.

\textbf{Linear SVC} matches RBF SVC on test accuracy (0.9825) with a very small $C{=}0.01$, indicating that strong regularization of the linear margin yields a well-generalizing model.

\textbf{RBF SVC} achieves the best ROC--AUC (0.9977) with $C{=}10$ and $\gamma{=}0.01$.
The moderate $\gamma$ keeps the kernel bandwidth wide enough to avoid overfitting, while the larger $C$ allows a tighter margin fit.

\subsection{Linear vs.\ Nonlinear Boundaries}

The comparable performance of linear models (Logistic, LinearSVC) to RBF SVC suggests the two classes are \emph{nearly linearly separable} in the 30-dimensional feature space.
The slight AUC advantage of RBF SVC indicates the nonlinear kernel captures subtle boundary curvature that linear models miss, though this translates to identical hard-label accuracy (0.9825).

\subsection{Regularization Insights}

All models favor moderate-to-strong regularization:
Logistic $C{=}0.1$, LinearSVC $C{=}0.01$, RBF SVC $C{=}10$ (moderate given the wider kernel).
This is expected: with only 455 training samples across 30 features, regularization prevents overfitting to noise.

%----------------------------------------------------------------------
\section{Reproducibility}
%----------------------------------------------------------------------

\begin{code}
SEED      = 42            # Fixes all random states
TEST_SIZE = 0.20          # 80/20 stratified split
N_FOLDS   = 5             # StratifiedKFold(shuffle=True, seed=42)
# Pipeline: StandardScaler -> model (fit inside each CV fold)
# GridSearchCV: scoring="accuracy", refit=True, n_jobs=-1
\end{code}

Three sources of randomness are fixed:
(1)~\texttt{random\_state=42} in \texttt{train\_test\_split} ensures identical partitions;
(2)~\texttt{random\_state=42} in \texttt{StratifiedKFold} ensures identical folds;
(3)~\texttt{random\_state=42} in each model ensures deterministic optimization.

%----------------------------------------------------------------------
\section{Appendix: Full Source Code}
%----------------------------------------------------------------------

\begin{code}
"""
PML Assignment 3 - Logistic Regression vs SVM
Breast Cancer Wisconsin Data
"""
import numpy as np, pandas as pd
from sklearn.datasets import load_breast_cancer
from sklearn.model_selection import (train_test_split,
    GridSearchCV, StratifiedKFold)
from sklearn.preprocessing import StandardScaler
from sklearn.pipeline import Pipeline
from sklearn.linear_model import LogisticRegression
from sklearn.svm import LinearSVC, SVC
from sklearn.calibration import CalibratedClassifierCV
from sklearn.metrics import (accuracy_score, confusion_matrix,
    roc_auc_score)
import warnings; warnings.filterwarnings("ignore")

SEED, TEST_SIZE, N_FOLDS = 42, 0.20, 5
C_GRID     = [0.001, 0.01, 0.1, 1, 10, 100]
GAMMA_GRID = [0.001, 0.01, 0.1, 1]

data = load_breast_cancer()
X, y = data.data, data.target

# 80/20 stratified split; test set used once at the end
X_train, X_test, y_train, y_test = train_test_split(
    X, y, test_size=TEST_SIZE, random_state=SEED, stratify=y)

cv = StratifiedKFold(n_splits=N_FOLDS, shuffle=True,
                     random_state=SEED)

def tune(name, pipe, grid):
    """Run GridSearchCV and print best result."""
    gs = GridSearchCV(pipe, grid, cv=cv, scoring="accuracy",
                      refit=True, n_jobs=-1)
    gs.fit(X_train, y_train)
    print(f"{name}: best_cv={gs.best_score_:.4f} "
          f"params={gs.best_params_}")
    return gs

# Model 1: Logistic Regression
gs_lr = tune("LR", Pipeline([
    ("scaler", StandardScaler()),
    ("clf", LogisticRegression(max_iter=10000,
                               random_state=SEED))
]), {"clf__C": C_GRID})

# Model 2: Linear SVC (calibrated for predict_proba)
gs_lsvc = tune("LSVC", Pipeline([
    ("scaler", StandardScaler()),
    ("clf", CalibratedClassifierCV(
        LinearSVC(max_iter=50000, dual="auto",
                  random_state=SEED), cv=3))
]), {"clf__estimator__C": C_GRID})

# Model 3: RBF SVC
gs_rbf = tune("RBF", Pipeline([
    ("scaler", StandardScaler()),
    ("clf", SVC(kernel="rbf", probability=True,
                random_state=SEED))
]), {"clf__C": [0.01,0.1,1,10,100],
     "clf__gamma": GAMMA_GRID})

# Final test evaluation (single pass)
for name, gs in [("LR",gs_lr), ("LSVC",gs_lsvc),
                 ("RBF",gs_rbf)]:
    yp = gs.predict(X_test)
    yprob = gs.predict_proba(X_test)[:,1]
    print(f"{name}: acc={accuracy_score(y_test,yp):.4f}"
          f" auc={roc_auc_score(y_test,yprob):.4f}")
    print(confusion_matrix(y_test, yp))
\end{code}

\end{document}
